This chapter addresses the \emph{control-abstraction} dimension of smart space control (\Cref{intro:sec:background}) ---\emph{what does the underlying control primitive communicate?}---and applies across all deployment types in \Cref{intro:tab:taxonomy}. We present \textbf{RASC}~\cite{rasc} (Request-Ack-Start-Complete) and its implementation \sysname, a new abstraction and system for IoT device control that addresses a fundamental gap in today's smart home platforms: the inability to know whether a physical action has actually taken effect. IoT platforms today rely on RPC, a fire-and-forget primitive that gives no feedback on action progress or failure---leaving both users and smart home developers without the observability needed to build reliable routines, reason about device state, or diagnose failures. \sysname{} exposes the full lifecycle of every action and introduces DAG-aware scheduling that enables automations to express causal dependencies, bringing observability and programmability to IoT settings.

This chapter is organized as follows. \Cref{rasc:sec:intro} introduces the problem and presents motivating examples. \Cref{rasc:sec:model} defines the system model. \Cref{rasc:sec:properties} states our design goals. \Cref{rasc:sec:rasc} describes the RASC abstraction and our adaptive polling strategy for action progress detection. \Cref{rasc:sec:casestudies} presents the programmability layer built atop RASC, including causality abstractions, dynamic scheduling, and our DAG-TL scheduler and STF and RV reschedulers. \Cref{rasc:sec:impl} describes our Home Assistant integration. \Cref{rasc:sec:eval} reports our experimental evaluation. \Cref{rasc:sec:discussion} discusses design considerations. \Cref{rasc:sec:related-work} surveys related work and \Cref{rasc:sec:conclusion} concludes.
